\section{Introduction}
\label{sec:intro}

By the Davis--Putnam--Robinson--Matiyasevich theorem \cite{matiyasevich1970},
every recursively enumerable set is Diophantine. The set of primes is therefore
Diophantine, and in 1976 Jones, Sato, Wada and Wiens made this concrete: they
wrote down a system of $14$ polynomial equations in $25$ unknowns and a
parameter whose solvability characterizes the primes, and packaged it into a
single polynomial whose positive values are exactly the primes
\cite{jsww1976}.

\subsection{The object of study is a frontier, not a number}

A polynomial $Q$ whose positive values over nonnegative integer arguments are
exactly $S$ is a \emph{generator} of $S$. Every generator carries two numbers:
its number of variables $\nu$ and its degree $\delta$. These two are in
tension. Given a representation
\[
  n\in S \iff \exists x_1,\dots,x_\nu\in\N:\; P_1=\dots=P_m=0 ,
\]
one obtains a generator by the classical device
$Q=W\cdot\bigl(1-\sum_i P_i^2\bigr)$, so that
$\deg Q = 1 + 2\max_i \deg P_i$; and the two operations available on such a
system move $(\nu,\delta)$ in opposite directions:

\begin{itemize}[nosep]
  \item \emph{flattening}---introducing a new unknown for a subexpression, with
        its defining equation---lowers $\max_i\deg P_i$ but \emph{adds}
        variables;
  \item \emph{elimination}---substituting away an unknown that some equation
        determines linearly---removes a variable but \emph{raises} the degree
        wherever that unknown occurred.
\end{itemize}

Consequently there is no single ``record''. What the literature describes is a
Pareto frontier of incomparable pairs, and a new pair is a contribution exactly
when it is not dominated by a known one.

\subsection{Announced versus exhibited}

The pairs cited for the primes below $26$ variables or below degree $25$ are,
with one exception, \emph{announced}: they are asserted to exist, sometimes with
a sketch of the construction, but the system itself was never written down.
Jones, Sato, Wada and Wiens themselves state $(42,5)$ and $(19,29)$ in a single
sentence on p.~449 of \cite{jsww1976}, and $(12,13697)$ follows from their
Theorem~2 without being exhibited; they write plainly that ``no literature is
available'' for the latter. The one exception at the low-variable end is
Matiyasevich's generator in $10$ variables \cite{matiyasevich1981}, which is
genuine and has since been formalized in Mizar \cite{pak2022ten}.

This distinction is not pedantry. An announced pair cannot be checked, cannot be
formalized, and cannot be used as the starting point for further work. The
present paper is an attempt to convert two announced regions of the frontier
into exhibited, machine-checked ones, and---in the course of doing so---to
correct what the announcements say.

\subsection{Contributions}

\begin{enumerate}[label=(C\arabic*),leftmargin=*]
  \item \textbf{A generator of the primes in $44$ variables of degree $5$}
        (Section~\ref{sec:deg5}), exhibited in full as a system of $32$
        equations of degree at most $2$ in $44$ variables. We are not aware of
        any previously exhibited prime generator of degree~$5$.

  \item \textbf{A generator in $21$ variables of degree $25$}
        (Section~\ref{sec:var21}), five fewer variables than the $(26,25)$ of
        \cite{jsww1976}, which is the only exhibited generator of the primes at
        that degree. An intermediate point $(23,25)$ is obtained with no
        auxiliary number theory at all, and is therefore the more robust of the
        two.

  \item \textbf{Formal verification in Lean~4} (Section~\ref{sec:lean}) of both
        constructions: for each, the equisatisfiability of the transformed
        system with the published one is a machine-checked theorem, using only
        the Lean kernel---no Mathlib, no \texttt{sorry}, and no axioms beyond
        the three the kernel itself uses. The Pell bound that licenses two of
        the five eliminations is proved from scratch rather than cited.

  \item \textbf{Soundness of the fourteen substitutions in the quotient-method
        system} (Section~\ref{sec:sec3}). Jones, Sato, Wada and Wiens dispose of
        these with the sentence ``the unknowns \dots\ eliminate from
        (I)--(XXI) by substitution''; over $\N$ each substitution requires the
        substituted value to be nonnegative, and eight of the fourteen do not
        follow from the shape of the equation. We prove and formalize all eight.

  \item \textbf{Three corrections to the literature}
        (Section~\ref{sec:correc}), each checkable against primary sources: the
        announced $(42,5)$ does not follow from the substitution method cited
        for it; and two circulating figures conflate a prime-representing
        polynomial with a universal Diophantine pair.
\end{enumerate}

\subsection{What we do not claim}

None of the pairs reported here is a proved minimum. Each is a constructed upper
bound, obtained by a search whose optimality statements are relative to the
search's own encoding. Section~\ref{sec:limits} is devoted to this and to the
history of figures we had to withdraw; we recommend reading it before the
constructions.
